\documentclass[10pt,a4paper]{amsart}
\usepackage[margin=19mm]{geometry}
\usepackage{amsmath,amssymb,mathtools}
\usepackage[scr=rsfs,cal=euler]{mathalfa}
\usepackage{xfrac}
\usepackage[unicode,hidelinks,bookmarks=false]{hyperref}
\newcommand{\ie}{i.e.\ }
\newcommand{\cf}{cf.\ }
\newcommand{\resp}{respectively\ }
\newcommand{\Subsection}{\subsection}
\providecommand{\nobreakdash}{\nobreak\hbox{-}}
\setlength{\emergencystretch}{2em}
\multlinegap0pt
\usepackage{mathtools}
\usepackage{amssymb}
\usepackage[shortlabels]{enumitem}
\usepackage{xcolor}
\usepackage{tikz}
\newtheorem{lemma}{Lemma}
\newtheorem{corollary}{Corollary}
\newtheorem{proposition}{Proposition}
\newcommand{\GL}{\operatorname{GL}}
\newcommand{\Q}{\mathbb{Q}}
\newcommand{\Supp}{\operatorname{Supp}}
\newcommand{\Z}{\mathbb{Z}}
\DeclareMathOperator{\im}{im}
\newcommand{\ord}{\operatorname{ord}}
\title{On the Finiteness of Isolated $j$-invariants for $X_1(N)$}
\author{Abbey Bourdon}
\date{Source: arXiv:2608.19354v1 (CC BY 4.0)}
\begin{document}
\maketitle

\noindent\emph{Source and licence.} On the Finiteness of Isolated $j$-invariants for $X_1(N)$. Version arXiv:2608.19354v1 (CC BY 4.0), distributed by arXiv under the Creative Commons Attribution 4.0 International licence, http://creativecommons.org/licenses/by/4.0/, as recorded in the arXiv licence metadata for this exact version and shown on the abstract page https://arxiv.org/abs/2608.19354v1. Full licence notice: \texttt{licenses/arxiv-2608.19354-CC-BY-4.0.txt}.\par\smallskip

\noindent\textit{Editorial scope.} Counted unit: the article's proposition LowerBoundProp (source0.tex lines 490--494). The article's complete original proof of it is delivered with it (source0.tex lines 496--506, one \texttt{proof} environment of 11 lines). No mathematics is written, repaired or re-derived: the statement and the proof are the article's own lines, in the article's own notation, and no retained line is substituted or re-flowed.  Retained uncounted as the source-local context the proof needs: lines 329--337; lines 443--447.  Nothing was omitted from any retained span, so the package carries no omission record.  No retained line contains a citation, so the wrapper carries no bibliography and no bibliography entry.  No retained line contains a figure environment, an image-inclusion command or drawing code, so no image is delivered and the package declares no asset dependency.  Editorial formatting only, in addition: an AMS \texttt{amsart} preamble that restates the article's own declarations and macros used by the retained lines, namely the environments \texttt{lemma}, \texttt{corollary}, \texttt{proposition} on separate counters, together with the standard packages those lines need.  The retained environments print in this document as Lemma 1, Corollary 1, Proposition 1 in order of appearance, whereas the article prints the same items with the numbering of its own full document.  The complete native source of the article as distributed in the arXiv e-print for this exact version is bundled unchanged as \texttt{sources/208079/source0.tex}.
\begin{lemma}\label{Finite_Degree_Bounds}
Fix $d \in \Z^+$, and let $\mathcal{E}$ be a collection of non-CM elliptic curves such that each $E\in \mathcal{E}$ is defined over $\Q(j(E))$ with $[\Q(j(E)):\Q]=d$. Let $S$ be a finite list of primes. Then there exists a constant $C=C(d)$ such that for any $x\in X_1(N)$ associated to $E \in \mathcal{E}$ with
\[
\Supp(N) \subseteq S \cup \{p: \im \rho_{E, p^{\infty}}=\GL_2(\Z_{p}) \}
\] one has
   \[
    \deg(x)\geq \frac{d}{C} \cdot \deg(X_1(N) \rightarrow X_1(1)).
    \]
\end{lemma}

\begin{corollary}\label{IntroCor}
Let $x\in X_1(N)$ be a non-CM point with $j(x)\in\Q$. Suppose $q \mid N$ for some prime $q > 37$, and set $\ord_{q}(N)=b$. Then \[
\deg(x)\geq\frac{1}{6}\deg(f(x))\deg(f),
\] where $f:X_1(N) \rightarrow X_1(N/q^b)$ is the natural map.    
\end{corollary}

\begin{proposition}\label{LowerBoundProp}
   There exists an absolute constant $C$ such that for any non-CM elliptic curve $E/\Q$ and any $x=[E,P]\in X_1(N)$ we have
    \[
    \deg(x) \geq C \cdot N^2\cdot \prod_{p \mid N}\frac{1}{6}\left(1-\frac{1}{p^2}\right) \]
\end{proposition}

\begin{proof}
   Let $E/\Q$ be a non-CM elliptic curve. Replacing $E$ by a twist if necessary, we may assume $-I \in \im \rho_{E,N}$. Set $S \coloneqq \{\text{primes} \leq 37\}$. Let $N=n_1\cdot n_2$, where $\Supp(n_1) \subseteq S$ and $\Supp(n_2) \cap S=\varnothing$. By repeatedly apply Corollary \ref{IntroCor}, we see 
        \[
        [\Q(P):\Q(n_2P)]\geq n_2^2\prod_{p \mid n_2}\frac{1}{6}\left(1-\frac{1}{p^2}\right).     \]
        By Lemma \ref{Finite_Degree_Bounds}, there is a constant $C_0$ which does not depend on $E$ such that
 \[
        [\Q(n_2P):\Q] \geq \frac{1}{C_0} \deg(X_1(n_1) \rightarrow X_1(1)) \geq \frac{1}{C_0} \cdot\frac{1}{2} \cdot n_1^2\cdot \prod_{p \mid n_1}\frac{1}{6}\left(1-\frac{1}{p^2}\right).
        \]
        Since $\deg(x)$ is at least $\frac{1}{2}\cdot [\Q(P):\Q] =\frac{1}{2}\cdot[\Q(P):\Q(n_2P)]\cdot[\Q(n_2P):\Q]$, the result follows. \qedhere
  
\end{proof}

\end{document}
